% GENERATED FILE. Edit canonical lessons, then run npm run generate:books.
% canonical-source-hash: fnv1a64:88a60807138305b9
% canonical-lessons: GU-C500-varso, GU-C500-vidhva, GU-C500-anath, GU-C500-dattak, GU-C500-tyajvu, GU-R500-review-494-500, GU-R500-review-final-500

\chapter{What They Left Us}
\label{ch:gu-what-they-left-us}
\hlchaptermodality{500}

\begin{chapteropening}
\textbf{By the end of this chapter:} \emph{I can talk about what a family passes on.}

Complete the last lesson of chapter 500: I can talk about what a family passes on.
\end{chapteropening}

\section[\texorpdfstring{\gu{વારસો} (vārso)}{વારસો (vārso)}]{\gu{વારસો} (vārso) — an inheritance; a heritage}

\label{lesson:GU-C500-varso}



\begin{quote}
Before the new one: say the Gujarati for a prayer meeting after a death, then the Gujarati for to offer with respect.
\end{quote}

\subsection*{You'll want to know: \gu{વારસો}}
\textbf{\gu{વારસો}} — \emph{vārso} — \textquotedblleft{}an inheritance; a heritage\textquotedblright{}.

\begin{quote}
\gu{દાદાએ} \gu{અમને} \gu{પુસ્તકોનો} \gu{વારસો} \gu{આપ્યો}.
\end{quote}

\begin{quote}
\emph{dādāe amne pustakono vārso āpyo.}
\end{quote}

\begin{quote}
Grandfather left us a heritage of books.
\end{quote}

A masculine noun.

\begin{cousinweb}
\textbf{\gu{વારસો}} (\emph{vārso}) is usually traced to Arabic \emph{wāriṯ}, \textquotedblleft{}heir\textquotedblright{}, which reached Gujarati through Persian \emph{wāris}.
\end{cousinweb}

\begin{grammarlens}[title={What you've built}]
One more word for describing what you lived and felt.
\end{grammarlens}

\subsection*{Your turn}
\begin{itemize}
\raggedright
  \item \emph{Say it:} \emph{vārso}
  \item \emph{Say it:} \emph{vārso}, once more
  \item \emph{Say it:} \emph{dādāe amne pustakono vārso āpyo.} — Grandfather left us a heritage of books.
  \item \emph{Recall:} say the Gujarati for a prayer meeting after a death, then the Gujarati for to offer with respect, then say \emph{vārso} again
\end{itemize}

\begin{tcolorbox}[breakable,colback=teal!4,colframe=teal!35!black,title={Before you move on}]
What does \gu{વારસો} mean? (\textquotedblleft{}An inheritance; a heritage\textquotedblright{}.) Say it once more.
\end{tcolorbox}
\section[\texorpdfstring{\gu{વિધવા} (vidhvā)}{વિધવા (vidhvā)}]{\gu{વિધવા} (vidhvā) — a widow}

\label{lesson:GU-C500-vidhva}



\begin{quote}
Before the new one: say the Gujarati for to offer with respect, then the Gujarati for an inheritance.
\end{quote}

\subsection*{You'll want to know: \gu{વિધવા}}
\textbf{\gu{વિધવા}} — \emph{vidhvā} — \textquotedblleft{}a widow\textquotedblright{}.

\begin{quote}
\gu{પતિ} \gu{ગુજરી} \gu{ગયા} \gu{પછી} \gu{મારી} \gu{કાકી} \gu{વિધવા} \gu{થઈ}.
\end{quote}

\begin{quote}
\emph{pati gujrī gayā pachī mārī kākī vidhvā thaī.}
\end{quote}

\begin{quote}
My aunt was widowed when her husband passed away.
\end{quote}

A feminine noun.

\begin{cousinweb}
\textbf{\gu{વિધવા}} (\emph{vidhvā}) is a tatsama, Sanskrit \emph{vidhavā}, \textquotedblleft{}a widow\textquotedblright{}, as Monier-Williams records it.
\end{cousinweb}

\begin{grammarlens}[title={What you've built}]
One more word for describing what you lived and felt.
\end{grammarlens}

\subsection*{Your turn}
\begin{itemize}
\raggedright
  \item \emph{Say it:} \emph{vidhvā}
  \item \emph{Say it:} \emph{vidhvā}, once more
  \item \emph{Say it:} \emph{pati gujrī gayā pachī mārī kākī vidhvā thaī.} — My aunt was widowed when her husband passed away.
  \item \emph{Recall:} say the Gujarati for to offer with respect, then the Gujarati for an inheritance, then say \emph{vidhvā} again
\end{itemize}

\begin{tcolorbox}[breakable,colback=teal!4,colframe=teal!35!black,title={Before you move on}]
What does \gu{વિધવા} mean? (\textquotedblleft{}A widow\textquotedblright{}.) Say it once more.
\end{tcolorbox}
\section[\texorpdfstring{\gu{અનાથ} (anāth)}{અનાથ (anāth)}]{\gu{અનાથ} (anāth) — an orphan}

\label{lesson:GU-C500-anath}



\begin{quote}
Before the new one: say the Gujarati for an inheritance, then the Gujarati for a widow.
\end{quote}

\subsection*{You'll want to know: \gu{અનાથ}}
\textbf{\gu{અનાથ}} — \emph{anāth} — \textquotedblleft{}an orphan\textquotedblright{}.

\begin{quote}
\gu{અમારા} \gu{ગામમાં} \gu{અનાથ} \gu{બાળકો} \gu{માટે} \gu{એક} \gu{ઘર} \gu{છે}.
\end{quote}

\begin{quote}
\emph{amārā gāmmā̃ anāth bāḷko māṭe ek ghar che.}
\end{quote}

\begin{quote}
In our village there is a home for orphaned children.
\end{quote}

A masculine or feminine noun, by the person meant.

\begin{cousinweb}
\textbf{\gu{અનાથ}} (\emph{anāth}) is a tatsama, Sanskrit \emph{anātha}, \textquotedblleft{}having no master or protector\textquotedblright{}, as Monier-Williams records it.
\end{cousinweb}

\begin{grammarlens}[title={What you've built}]
One more word for describing what you lived and felt.
\end{grammarlens}

\subsection*{Your turn}
\begin{itemize}
\raggedright
  \item \emph{Say it:} \emph{anāth}
  \item \emph{Say it:} \emph{anāth}, once more
  \item \emph{Say it:} \emph{amārā gāmmā̃ anāth bāḷko māṭe ek ghar che.} — In our village there is a home for orphaned children.
  \item \emph{Recall:} say the Gujarati for an inheritance, then the Gujarati for a widow, then say \emph{anāth} again
\end{itemize}

\begin{tcolorbox}[breakable,colback=teal!4,colframe=teal!35!black,title={Before you move on}]
What does \gu{અનાથ} mean? (\textquotedblleft{}An orphan\textquotedblright{}.) Say it once more.
\end{tcolorbox}
\section[\texorpdfstring{\gu{દત્તક} (dattak)}{દત્તક (dattak)}]{\gu{દત્તક} (dattak) — adopted}

\label{lesson:GU-C500-dattak}



\begin{quote}
Before the new one: say the Gujarati for a widow, then the Gujarati for an orphan.
\end{quote}

\subsection*{You'll want to know: \gu{દત્તક}}
\textbf{\gu{દત્તક}} — \emph{dattak} — \textquotedblleft{}adopted\textquotedblright{}.

\begin{quote}
\gu{તેમણે} \gu{એક} \gu{અનાથ} \gu{બાળકને} \gu{દત્તક} \gu{લીધું}.
\end{quote}

\begin{quote}
\emph{temṇe ek anāth bāḷakne dattak līdhũ.}
\end{quote}

\begin{quote}
They adopted an orphaned child.
\end{quote}

\begin{cousinweb}
\textbf{\gu{દત્તક}} (\emph{dattak}) is a tatsama, Sanskrit \emph{dattaka}, which Monier-Williams gives, with \emph{putra}, \textquotedblleft{}son\textquotedblright{}, as an adopted son.
\end{cousinweb}

\begin{grammarlens}[title={What you've built}]
One more word for describing what you lived and felt.
\end{grammarlens}

\subsection*{Your turn}
\begin{itemize}
\raggedright
  \item \emph{Say it:} \emph{dattak}
  \item \emph{Say it:} \emph{dattak}, once more
  \item \emph{Say it:} \emph{temṇe ek anāth bāḷakne dattak līdhũ.} — They adopted an orphaned child.
  \item \emph{Recall:} say the Gujarati for a widow, then the Gujarati for an orphan, then say \emph{dattak} again
\end{itemize}

\begin{tcolorbox}[breakable,colback=teal!4,colframe=teal!35!black,title={Before you move on}]
What does \gu{દત્તક} mean? (\textquotedblleft{}Adopted\textquotedblright{}.) Say it once more.
\end{tcolorbox}
\section[\texorpdfstring{\gu{ત્યજવું} (tyajvũ)}{ત્યજવું (tyajvũ)}]{\gu{ત્યજવું} (tyajvũ) — to abandon; to give up}

\label{lesson:GU-C500-tyajvu}



\begin{quote}
Before the new one: say the Gujarati for an orphan, then the Gujarati for adopted.
\end{quote}

\subsection*{You'll want to know: \gu{ત્યજવું}}
\textbf{\gu{ત્યજવું}} — \emph{tyajvũ} — \textquotedblleft{}to abandon; to give up\textquotedblright{}.

\begin{quote}
\gu{તેણે} \gu{ગુસ્સો} \gu{ત્યજી} \gu{દીધો}.
\end{quote}

\begin{quote}
\emph{teṇe gusso tyajī dīdho.}
\end{quote}

\begin{quote}
He gave up his anger.
\end{quote}

Past, with \gu{તેણે} (he or she …): \gu{તેણે} \gu{ઘર} \gu{ત્યજ્યું}, \emph{teṇe ghar tyajyũ}.

\begin{cousinweb}
\textbf{\gu{ત્યજવું}} (\emph{tyajvũ}) goes back to Sanskrit \emph{tyajati}, \textquotedblleft{}leaves, abandons\textquotedblright{}, which Monier-Williams records under \emph{tyaj}.
\end{cousinweb}

\begin{grammarlens}[title={What you've built}]
One more word for describing what you lived and felt.
\end{grammarlens}

\subsection*{Your turn}
\begin{itemize}
\raggedright
  \item \emph{Say it:} \emph{tyajvũ}
  \item \emph{Say it:} \emph{tyajvũ}, once more
  \item \emph{Say it:} \emph{teṇe gusso tyajī dīdho.} — He gave up his anger.
  \item \emph{Recall:} say the Gujarati for an orphan, then the Gujarati for adopted, then say \emph{tyajvũ} again
\end{itemize}

\begin{tcolorbox}[breakable,colback=teal!4,colframe=teal!35!black,title={Before you move on}]
What does \gu{ત્યજવું} mean? (\textquotedblleft{}To abandon; to give up\textquotedblright{}.) Say it once more.
\end{tcolorbox}
\section[(dialogue)]{Looking back: from Growing Up, Growing Old to What They Left Us}

\label{lesson:GU-R500-review-494-500}



\begin{quote}
Say the Gujarati for childhood and to grow up.
\end{quote}

\subsection*{Your turn}
Say these aloud:

\begin{itemize}
\raggedright
  \item childhood, to grow up, a teenage boy, elderly, an old people's home
  \item pregnant, childbirth, a baby-shower ceremony, to come into the world, to give birth to
  \item to suckle, to rock a cradle, a naming ceremony, to welcome with rites, a baby's nappy
  \item a wedding invitation card, the turmeric rite, the groom's wedding party, the groom's procession, to marry off
  \item a farewell, a bride, unmarried, a dowry, a divorce
  \item a death, a passing away, to pass away, a prayer meeting after a death, to offer with respect
  \item an inheritance, a widow, an orphan, adopted, to abandon
\end{itemize}

\begin{tcolorbox}[breakable,colback=teal!4,colframe=teal!35!black,title={Before you move on}]
Childhood? (\textbf{\gu{બાળપણ}}, \emph{bāḷpaṇ}.) To grow up? (\textbf{\gu{ઉછરવું}}, \emph{ucharvũ}.) A teenage boy? (\textbf{\gu{કિશોર}}, \emph{kiśor}.) Elderly? (\textbf{\gu{વૃદ્ધ}}, \emph{vṛddh}.) An old people's home? (\textbf{\gu{વૃદ્ધાશ્રમ}}, \emph{vṛddhāśram}.) Pregnant? (\textbf{\gu{સગર્ભા}}, \emph{sagarbhā}.) Childbirth? (\textbf{\gu{પ્રસૂતિ}}, \emph{prasūti}.) A baby-shower ceremony? (\textbf{\gu{સીમંત}}, \emph{sīmant}.) To come into the world? (\textbf{\gu{અવતરવું}}, \emph{avtarvũ}.) To give birth to? (\textbf{\gu{જણવું}}, \emph{jaṇvũ}.) To suckle? (\textbf{\gu{ધાવવું}}, \emph{dhāvvũ}.) To rock a cradle? (\textbf{\gu{હીંચોળવું}}, \emph{hī̃coḷvũ}.) A naming ceremony? (\textbf{\gu{નામકરણ}}, \emph{nāmkaraṇ}.) To welcome with rites? (\textbf{\gu{વધાવવું}}, \emph{vadhāvvũ}.) A baby's nappy? (\textbf{\gu{બાળોતિયું}}, \emph{bāḷotiyũ}.) A wedding invitation card? (\textbf{\gu{કંકોતરી}}, \emph{kankotrī}.) The turmeric rite? (\textbf{\gu{પીઠી}}, \emph{pīṭhī}.) The groom's wedding party? (\textbf{\gu{જાન}}, \emph{jān}.) The groom's procession? (\textbf{\gu{વરઘોડો}}, \emph{varghoḍo}.) To marry off? (\textbf{\gu{પરણાવવું}}, \emph{parṇāvvũ}.) A farewell? (\textbf{\gu{વિદાય}}, \emph{vidāy}.) A bride? (\textbf{\gu{કન્યા}}, \emph{kanyā}.) Unmarried? (\textbf{\gu{કુંવારું}}, \emph{kunvārũ}.) A dowry? (\textbf{\gu{દહેજ}}, \emph{dahej}.) A divorce? (\textbf{\gu{છૂટાછેડા}}, \emph{chūṭācheḍā}.) A death? (\textbf{\gu{મરણ}}, \emph{maraṇ}.) A passing away? (\textbf{\gu{અવસાન}}, \emph{avsān}.) To pass away? (\textbf{\gu{ગુજરવું}}, \emph{gujarvũ}.) A prayer meeting after a death? (\textbf{\gu{બેસણું}}, \emph{besṇũ}.) To offer with respect? (\textbf{\gu{અર્પવું}}, \emph{arpavũ}.) An inheritance? (\textbf{\gu{વારસો}}, \emph{vārso}.) A widow? (\textbf{\gu{વિધવા}}, \emph{vidhvā}.) An orphan? (\textbf{\gu{અનાથ}}, \emph{anāth}.) Adopted? (\textbf{\gu{દત્તક}}, \emph{dattak}.) To abandon? (\textbf{\gu{ત્યજવું}}, \emph{tyajvũ}.)
\end{tcolorbox}
\section[(dialogue)]{What They Left Us — once more}

\label{lesson:GU-R500-review-final-500}



\begin{quote}
Say the Gujarati for an inheritance and a widow.
\end{quote}

\subsection*{Your turn}
\begin{itemize}
\raggedright
  \item \emph{Say it:} an inheritance, a widow, an orphan, adopted, to abandon
\end{itemize}

\begin{tcolorbox}[breakable,colback=teal!4,colframe=teal!35!black,title={Before you move on}]
An inheritance? (\textbf{\gu{વારસો}}, \emph{vārso}.) A widow? (\textbf{\gu{વિધવા}}, \emph{vidhvā}.) An orphan? (\textbf{\gu{અનાથ}}, \emph{anāth}.) Adopted? (\textbf{\gu{દત્તક}}, \emph{dattak}.) To abandon? (\textbf{\gu{ત્યજવું}}, \emph{tyajvũ}.)
\end{tcolorbox}
